\chapter{Dynamic Programming: DP on LIS}
\chapterepigraph{The longest increasing subsequence asks one question
at every position: looking back at everything already placed, what is
the best chain I could be extending right now? Once that single
question is answered efficiently, an entire family of problems --
divisibility chains, string chains, bitonic sequences, counting -- falls
out as the same question asked with a different notion of ``extends.''}

\section{Ending-Here DP: A New Way to Index the Table}
\label{sec:dp-lis-intro}

Every problem in this chapter uses a state shape not yet seen in this
book: $\textit{dp}[i]$ represents the best answer for a subsequence (or
chain) \textbf{ending exactly at index $i$} -- not a prefix-based answer
like every earlier chapter's $f(i, \ldots)$, which represented the best
answer using \emph{only} the first $i$ elements. This ``ending-here''
indexing is what allows the final answer to be the maximum over
\emph{all} $\textit{dp}[i]$, rather than simply $\textit{dp}[n]$.

\begin{ideabox}[Why the Answer Lives in a Maximum Over the Whole Table, Not One Cell]
A longest increasing subsequence can end at \emph{any} position in the
array -- there is no reason to expect it ends at the last element. So,
unlike LCS or Knapsack-family problems where the final answer sits at
one specific corner of the table ($\textit{dp}[n][m]$), here every cell
$\textit{dp}[i]$ is independently a candidate answer, and we must scan
the entire filled table for the maximum at the end. This is the single
biggest structural difference to watch for in this chapter.
\end{ideabox}

\begin{patternbox}[How ``Ending-Here'' Generalizes Beyond Plain LIS]
\begin{itemize}
  \item \textbf{Longest Increasing Subsequence}: $\textit{dp}[i] = 1 +
        \max(\textit{dp}[j])$ over all $j<i$ with
        $\textit{arr}[j] < \textit{arr}[i]$.
  \item \textbf{Largest Divisible Subset (Section~\ref{sec:dp-largest-divisible-subset})}:
        identical shape, with the condition $\textit{arr}[j] <
        \textit{arr}[i]$ replaced by $\textit{arr}[i] \bmod
        \textit{arr}[j] = 0$ (after sorting).
  \item \textbf{Longest String Chain (Section~\ref{sec:dp-longest-string-chain})}:
        identical shape, with the condition replaced by ``$\textit{word}_j$
        becomes $\textit{word}_i$ after inserting exactly one
        character'' (after sorting by length).
  \item \textbf{Longest Bitonic Subsequence (Section~\ref{sec:dp-longest-bitonic-subsequence})}:
        run the ending-here LIS computation \textbf{twice} (once
        forward for increasing runs, once backward for decreasing
        runs) and combine the two tables at each index.
  \item \textbf{Number of LIS (Section~\ref{sec:dp-number-of-lis})}: pair
        $\textit{dp}[i]$ (the best length ending here) with a second
        table $\textit{count}[i]$ (how many subsequences of that best
        length end here), updated together.
\end{itemize}
Recognizing this single recurrence shape -- extend from any valid
earlier ``ending-here'' cell, under whatever notion of ``valid'' the
problem defines -- is the entire content of this chapter.
\end{patternbox}

\newpage

\section{Longest Increasing Subsequence}
\label{sec:dp-lis}

\problemstatement{Given an array $\textit{arr}$ of $n$ integers, find
the length of the \textbf{longest strictly increasing subsequence}
(elements in order, but not necessarily contiguous).}

\begin{patternbox}
\emph{``Longest increasing subsequence''} is the chapter's foundational
problem: at every index, the question is whether it can \textbf{extend}
some earlier subsequence -- which depends on what the \emph{previous
chosen element's value} was, not just which index we're at.
\end{patternbox}

\subsection*{Deriving the recurrence (pick/not-pick formulation)}

Let $f(\textit{ind}, \textit{prevIndex})$ be the length of the longest
increasing subsequence achievable from index $\textit{ind}$ onward,
given that the most recently \emph{included} element (if any) is at
index $\textit{prevIndex}$ (or $-1$ if nothing has been included yet).
At each index, either skip it ($f(\textit{ind}{+}1, \textit{prevIndex})$),
or include it -- only valid if $\textit{prevIndex}=-1$ or
$\textit{arr}[\textit{prevIndex}] < \textit{arr}[\textit{ind}]$ -- gaining
$1 + f(\textit{ind}{+}1, \textit{ind})$:
\[
  f(\textit{ind}, \textit{prevIndex}) = \max\Big(
  f(\textit{ind}{+}1, \textit{prevIndex}),\ \textit{canTake} \Big), \quad
  \textit{canTake} = 1 + f(\textit{ind}{+}1, \textit{ind})\ \text{(if valid)}.
\]

\subsection*{Approach 1: Recursion}

\begin{lstlisting}
#include <bits/stdc++.h>
using namespace std;

int lisRecursive(int ind, int prevIndex, vector<int>& arr) {
    int n = arr.size();
    if (ind == n) return 0;

    int notTake = lisRecursive(ind + 1, prevIndex, arr);
    int take = 0;
    if (prevIndex == -1 || arr[prevIndex] < arr[ind]) {
        take = 1 + lisRecursive(ind + 1, ind, arr);
    }
    return max(notTake, take);
}

int main() {
    vector<int> arr = {10, 9, 2, 5, 3, 7, 101, 18};
    cout << lisRecursive(0, -1, arr) << endl; // 4
    return 0;
}
\end{lstlisting}

\begin{complexitybox}[Approach 1 Complexity]
\textbf{Time.} $O(2^n)$.
\textbf{Space.} $O(n)$ recursion stack.
\end{complexitybox}

\subsection*{Approach 2: Memoization}

\begin{lstlisting}
#include <bits/stdc++.h>
using namespace std;

int lisMemo(int ind, int prevIndex, vector<int>& arr, vector<vector<int>>& memo) {
    int n = arr.size();
    if (ind == n) return 0;
    int prevShift = prevIndex + 1; // shift by 1 so -1 maps to 0
    if (memo[ind][prevShift] != -1) return memo[ind][prevShift];

    int notTake = lisMemo(ind + 1, prevIndex, arr, memo);
    int take = 0;
    if (prevIndex == -1 || arr[prevIndex] < arr[ind]) {
        take = 1 + lisMemo(ind + 1, ind, arr, memo);
    }
    return memo[ind][prevShift] = max(notTake, take);
}

int main() {
    vector<int> arr = {10, 9, 2, 5, 3, 7, 101, 18};
    int n = arr.size();
    vector<vector<int>> memo(n, vector<int>(n + 1, -1));
    cout << lisMemo(0, -1, arr, memo) << endl; // 4
    return 0;
}
\end{lstlisting}

\begin{complexitybox}[Approach 2 Complexity]
\textbf{Time.} $O(n^2)$.
\textbf{Space.} $O(n^2)$ memo $+$ $O(n)$ recursion stack.
\end{complexitybox}

\subsection*{Approach 3: Tabulation (Ending-Here Reformulation)}

\begin{ideabox}[Why Tabulation Uses a Cleaner, Different-Looking Table]
Directly tabulating the $(\textit{ind}, \textit{prevIndex})$ recursion
is awkward (the shifted-index bookkeeping is unpleasant to iterate over
bottom-up). The standard simplification: define
$\textit{dp}[i]$ = length of the LIS \emph{ending exactly at} index $i$.
This is mathematically equivalent information, just reorganized; it
gives a much cleaner loop and is how this problem is conventionally
tabulated.
\end{ideabox}

\[
  \textit{dp}[i] = 1 + \max_{j<i,\ \textit{arr}[j]<\textit{arr}[i]}
  \textit{dp}[j] \quad (\text{or } 1 \text{ if no such } j \text{ exists}).
\]

\begin{lstlisting}
#include <bits/stdc++.h>
using namespace std;

int lisTabulation(vector<int>& arr) {
    int n = arr.size();
    vector<int> dp(n, 1); // every element alone is an LIS of length 1

    int best = 1;
    for (int i = 0; i < n; i++) {
        for (int j = 0; j < i; j++) {
            if (arr[j] < arr[i]) {
                dp[i] = max(dp[i], 1 + dp[j]);
            }
        }
        best = max(best, dp[i]);
    }
    return best;
}

int main() {
    vector<int> arr = {10, 9, 2, 5, 3, 7, 101, 18};
    cout << lisTabulation(arr) << endl; // 4
    return 0;
}
\end{lstlisting}

\subsection*{Dry run of the tabulation table}

\dryrun{Take $\textit{arr} = [10,9,2,5,3,7,101,18]$.}

\begin{center}
\begin{tabular}{c|c|c|c|c|c|c|c|c}
$i$ & 0 & 1 & 2 & 3 & 4 & 5 & 6 & 7 \\ \hline
$\textit{arr}[i]$ & 10 & 9 & 2 & 5 & 3 & 7 & 101 & 18 \\ \hline
$\textit{dp}[i]$ & 1 & 1 & 1 & 2 & 2 & 3 & 4 & 4
\end{tabular}
\end{center}

$\textit{dp}[5]=3$ comes from extending the chain ending at index $3$
(value $5$, $\textit{dp}=2$): $2,5,7$. $\textit{dp}[6]=4$ extends that
same chain with $101$: $2,5,7,101$. The maximum over the whole table is
$4$.

\begin{complexitybox}[Approach 3 Complexity]
\textbf{Time.} $O(n^2)$ -- a nested loop comparing every pair.
\textbf{Space.} $O(n)$ for the \textit{dp} array -- already minimal;
this problem's natural tabulated form has no further space reduction
the way 2D-table problems in earlier chapters did.
\end{complexitybox}

\begin{ideabox}[A Faster $O(n \log n)$ Algorithm Exists]
Section~\ref{sec:dp-lis-binary-search} develops a completely different
$O(n \log n)$ technique for this same problem, using binary search
instead of dynamic programming's pairwise comparison -- worth knowing
as the production-quality solution once the $O(n^2)$ DP version here is
understood.
\end{ideabox}

\begin{takeawaybox}
``Ending-here'' tables answer their question with a \textbf{maximum
over the entire filled table}, not a single corner cell -- the first
new structural lesson of this chapter, and the one to watch for in
every section that follows.
\end{takeawaybox}

\section{Print the Longest Increasing Subsequence}
\label{sec:dp-print-lis}

\problemstatement{Given $\textit{arr}$, construct (not just measure the
length of) one actual longest increasing subsequence.}

\begin{patternbox}
Exactly as in the DP on Strings chapter's Print LCS section, reconstructing the actual
solution needs the full tabulation table kept around, plus a
\textbf{parent pointer} recorded at each index, tracking \emph{which
earlier index} extended into it -- since the table itself only stores
chain \emph{lengths}, not which specific chain produced them.
\end{patternbox}

\subsection*{Understanding the reconstruction}

Alongside $\textit{dp}[i]$ (the LIS length ending at $i$), maintain
$\textit{parent}[i]$: the index $j<i$ that $\textit{dp}[i]$ was
extended from (or $i$ itself, if $\textit{dp}[i]=1$ and nothing extends
it). After filling the table, find the index with the maximum
$\textit{dp}$ value, then follow $\textit{parent}$ pointers backward
until reaching an index whose parent is itself, collecting elements
along the way.

\subsection*{Algorithm}

\begin{enumerate}
  \item Initialize $\textit{dp}[i] \gets 1$, $\textit{parent}[i] \gets
        i$ for all $i$.
  \item For $i$ from $0$ to $n-1$: for $j$ from $0$ to $i-1$: if
        $\textit{arr}[j]<\textit{arr}[i]$ and $1+\textit{dp}[j] >
        \textit{dp}[i]$: update $\textit{dp}[i] \gets 1+\textit{dp}[j]$;
        $\textit{parent}[i] \gets j$.
  \item Find $\textit{lastIndex}$ maximizing $\textit{dp}[i]$.
  \item Walk backward from $\textit{lastIndex}$ via \textit{parent}
        until $\textit{parent}[i]=i$, collecting each $\textit{arr}[i]$
        visited; reverse the collected list.
\end{enumerate}

\subsection*{Dry run}

\dryrun{Using Section~\ref{sec:dp-lis}'s table for
$\textit{arr}=[10,9,2,5,3,7,101,18]$.}

$\textit{dp}=[1,1,1,2,2,3,4,4]$, with $\textit{parent}[3]=2$ (chain
$5$ extends $2$), $\textit{parent}[5]=3$ (chain $7$ extends $5$),
$\textit{parent}[6]=5$ (chain $101$ extends $7$). The maximum
$\textit{dp}$ value is $4$, first achieved at index $6$. Walking
backward: $6 \to 5 \to 3 \to 2$, collecting $101, 7, 5, 2$; reversing
gives $[2,5,7,101]$.

\subsection*{C++ implementation}

\begin{lstlisting}
#include <bits/stdc++.h>
using namespace std;

vector<int> printLIS(vector<int>& arr) {
    int n = arr.size();
    vector<int> dp(n, 1), parent(n);
    for (int i = 0; i < n; i++) parent[i] = i;

    int lastIndex = 0;
    for (int i = 0; i < n; i++) {
        for (int j = 0; j < i; j++) {
            if (arr[j] < arr[i] && 1 + dp[j] > dp[i]) {
                dp[i] = 1 + dp[j];
                parent[i] = j;
            }
        }
        if (dp[i] > dp[lastIndex]) lastIndex = i;
    }

    vector<int> result;
    int cur = lastIndex;
    while (parent[cur] != cur) {
        result.push_back(arr[cur]);
        cur = parent[cur];
    }
    result.push_back(arr[cur]);
    reverse(result.begin(), result.end());
    return result;
}

int main() {
    vector<int> arr = {10, 9, 2, 5, 3, 7, 101, 18};
    for (int x : printLIS(arr)) cout << x << " ";
    cout << endl; // 2 5 7 101
    return 0;
}
\end{lstlisting}

\begin{complexitybox}
\textbf{Time Complexity.} $O(n^2)$ to build the table, plus $O(n)$ for
the backward walk.
\textbf{Space Complexity.} $O(n)$ for the \textit{dp} and
\textit{parent} arrays.
\end{complexitybox}

\begin{takeawaybox}
``Print the actual chain'' problems on an ending-here table always need
a parent-pointer array recorded \emph{during} the forward fill, then a
backward walk from whichever index holds the global maximum -- the same
reconstruction idea as Print LCS, adapted to this chapter's different
table shape.
\end{takeawaybox}

\section{Longest Increasing Subsequence using Binary Search}
\label{sec:dp-lis-binary-search}

\problemstatement{Solve Section~\ref{sec:dp-lis} (find the length of the
LIS) in $O(n \log n)$ instead of $O(n^2)$.}

\begin{patternbox}
\emph{``Find the length of the LIS, faster than $O(n^2)$''} is solved
not by a better DP table, but by an entirely different technique:
maintain a single array $\textit{tails}$, where $\textit{tails}[k]$ is
the \textbf{smallest possible value} that can end an increasing
subsequence of length $k+1$, discovered so far -- and update it for
each new element using \textbf{binary search} (the Binary Search chapter's
lower bound) instead of scanning all previous elements.
\end{patternbox}

\subsection*{Understanding why \textit{tails} is always sorted}

The key invariant: $\textit{tails}$ is always kept sorted in increasing
order. This is not obvious at first, but follows by induction: a longer
increasing subsequence's smallest possible tail value can never be
smaller than a shorter one's smallest possible tail value (if it were,
that shorter prefix of the longer subsequence would itself be a
counterexample to the shorter subsequence's tail already being
minimal). Because $\textit{tails}$ stays sorted, we can binary search it.

\begin{ideabox}[Update Rule for Each New Element]
For each $x = \textit{arr}[i]$, in order: binary search $\textit{tails}$
for the leftmost position $\textit{pos}$ where $\textit{tails}[\textit{pos}]
\ge x$ (this is exactly the Binary Search chapter's lower bound technique).
If no such position exists ($x$ is larger than every current tail),
append $x$ (a new, longer subsequence is now possible). Otherwise,
replace $\textit{tails}[\textit{pos}] \gets x$ ($x$ is a \emph{smaller}
valid tail for a subsequence of that same length, strictly improving
future extension possibilities, without changing the \emph{count} of
distinct lengths discovered).
\end{ideabox}

\subsection*{Why replacing (not just appending) preserves correctness}

$\textit{tails}$'s length, at any point, equals the LIS length found so
far among the elements processed -- but the actual \emph{values} stored
are not necessarily a real subsequence from the array; they are simply
the best (smallest) tail achievable for each length. Replacing
$\textit{tails}[\textit{pos}]$ with a smaller value $x$ never decreases
$\textit{tails}$'s length (we are not removing entries, only improving
one in place), and it can only ever \emph{help} future elements extend a
chain of that length, since a smaller tail is easier for a later,
possibly-not-much-larger element to exceed. The final length of
$\textit{tails}$ is therefore exactly the LIS length.

\subsection*{Algorithm}

\begin{enumerate}
  \item Initialize an empty array \textit{tails}.
  \item For each $x$ in $\textit{arr}$: binary search for the lower
        bound of $x$ in \textit{tails}. If found at the end (no entry
        $\ge x$): append $x$. Else: overwrite that entry with $x$.
  \item Return $|\textit{tails}|$.
\end{enumerate}

\subsection*{Dry run}

\dryrun{Take $\textit{arr} = [10,9,2,5,3,7,101,18]$.}

\begin{itemize}
  \item $x=10$: \textit{tails} empty, append. $\textit{tails}=[10]$.
  \item $x=9$: lower bound of $9$ in $[10]$ is index $0$. Overwrite.
        $\textit{tails}=[9]$.
  \item $x=2$: lower bound of $2$ in $[9]$ is index $0$. Overwrite.
        $\textit{tails}=[2]$.
  \item $x=5$: lower bound of $5$ in $[2]$: none $\ge5$, append.
        $\textit{tails}=[2,5]$.
  \item $x=3$: lower bound of $3$ in $[2,5]$ is index $1$ (value $5$).
        Overwrite. $\textit{tails}=[2,3]$.
  \item $x=7$: lower bound of $7$ in $[2,3]$: none $\ge7$, append.
        $\textit{tails}=[2,3,7]$.
  \item $x=101$: append. $\textit{tails}=[2,3,7,101]$.
  \item $x=18$: lower bound of $18$ in $[2,3,7,101]$ is index $3$
        (value $101$). Overwrite. $\textit{tails}=[2,3,7,18]$.
\end{itemize}

Final length: $4$ -- matching Section~\ref{sec:dp-lis}'s answer
(note $\textit{tails}=[2,3,7,18]$ is not necessarily a real
subsequence of $\textit{arr}$, but its \emph{length} is correctly the
LIS length).

\subsection*{C++ implementation}

\begin{lstlisting}
#include <bits/stdc++.h>
using namespace std;

int lisBinarySearch(vector<int>& arr) {
    vector<int> tails;

    for (int x : arr) {
        int pos = lower_bound(tails.begin(), tails.end(), x) - tails.begin();
        if (pos == (int)tails.size()) {
            tails.push_back(x);
        } else {
            tails[pos] = x;
        }
    }
    return tails.size();
}

int main() {
    vector<int> arr = {10, 9, 2, 5, 3, 7, 101, 18};
    cout << lisBinarySearch(arr) << endl; // 4
    return 0;
}
\end{lstlisting}

\begin{complexitybox}
\textbf{Time Complexity.} Each of the $n$ elements performs one binary
search over \textit{tails} (size at most $n$), giving
\[
  O(n \log n).
\]
\textbf{Space Complexity.} $O(n)$ for the \textit{tails} array.
\end{complexitybox}

\begin{takeawaybox}
Not every problem's fastest solution is a bigger or cleverer DP table --
sometimes, as here, an entirely different paradigm (greedy maintenance
of a sorted auxiliary array, queried by binary search) beats DP's
polynomial bound outright. Keep this $O(n \log n)$ technique available
specifically for when a problem only needs the LIS \emph{length}; it
does not directly reconstruct the actual subsequence the way
Section~\ref{sec:dp-print-lis}'s parent-pointer approach does.
\end{takeawaybox}

\section{Largest Divisible Subset}
\label{sec:dp-largest-divisible-subset}

\problemstatement{Given an array of \textbf{distinct} positive
integers, find the largest subset such that for every pair of elements
$(a,b)$ in it, either $a \bmod b = 0$ or $b \bmod a = 0$.}

\begin{patternbox}
\emph{``Largest subset where every pair divides''} is
Section~\ref{sec:dp-lis-intro}'s ending-here template with the
condition $\textit{arr}[j] < \textit{arr}[i]$ replaced by
$\textit{arr}[i] \bmod \textit{arr}[j] = 0$ -- after first
\textbf{sorting} the array, exactly the same preparatory step LIS uses
implicitly (LIS's input is typically already read left to right; here
we must sort first to make the pairwise-divisibility condition reduce
to a simple left-to-right chain check).
\end{patternbox}

\subsection*{Why sorting first is essential}

After sorting, if $\textit{arr}[k]$ divides $\textit{arr}[j]$ and
$\textit{arr}[j]$ divides $\textit{arr}[i]$ (with $k<j<i$ in the sorted
order), then $\textit{arr}[k]$ \emph{automatically} divides
$\textit{arr}[i]$ too (divisibility is transitive). So checking only
\textbf{consecutive} chain links ($\textit{arr}[j] \mid
\textit{arr}[i]$ for the immediately preceding chosen element) is
sufficient to guarantee \emph{every pair} in the resulting chain
satisfies the divisibility condition -- exactly mirroring why LIS only
checks $\textit{arr}[j]<\textit{arr}[i]$ pairwise rather than every
triple.

\subsection*{The recurrence}

\[
  \textit{dp}[i] = 1 + \max_{j<i,\ \textit{arr}[i] \bmod
  \textit{arr}[j] = 0} \textit{dp}[j] \quad (\text{or } 1 \text{ if no
  such } j).
\]

\subsection*{Algorithm (with reconstruction)}

Exactly Section~\ref{sec:dp-print-lis}'s parent-pointer approach,
with the extension condition swapped to divisibility.

\begin{lstlisting}
#include <bits/stdc++.h>
using namespace std;

vector<int> largestDivisibleSubset(vector<int> arr) {
    int n = arr.size();
    sort(arr.begin(), arr.end());

    vector<int> dp(n, 1), parent(n);
    for (int i = 0; i < n; i++) parent[i] = i;

    int lastIndex = 0;
    for (int i = 0; i < n; i++) {
        for (int j = 0; j < i; j++) {
            if (arr[i] % arr[j] == 0 && 1 + dp[j] > dp[i]) {
                dp[i] = 1 + dp[j];
                parent[i] = j;
            }
        }
        if (dp[i] > dp[lastIndex]) lastIndex = i;
    }

    vector<int> result;
    int cur = lastIndex;
    while (parent[cur] != cur) {
        result.push_back(arr[cur]);
        cur = parent[cur];
    }
    result.push_back(arr[cur]);
    reverse(result.begin(), result.end());
    return result;
}

int main() {
    vector<int> arr = {1, 16, 7, 8, 4};
    for (int x : largestDivisibleSubset(arr)) cout << x << " ";
    cout << endl; // 1 4 8 16
    return 0;
}
\end{lstlisting}

\subsection*{Dry run}

\dryrun{Take $\textit{arr} = [1,16,7,8,4]$, sorted: $[1,4,7,8,16]$.}

$\textit{dp}=[1,2,1,3,4]$: $4 \bmod 1=0 \Rightarrow \textit{dp}[1]=2$;
$8 \bmod 4=0 \Rightarrow \textit{dp}[3]=3$ (extending the chain
$1,4,8$); $16 \bmod 8=0 \Rightarrow \textit{dp}[4]=4$ (extending
$1,4,8,16$). Note $7$ extends nothing (no earlier sorted element
divides it), so $\textit{dp}[2]=1$. The maximum is $4$, giving the
subset $[1,4,8,16]$.

\begin{complexitybox}
\textbf{Time Complexity.} $O(n \log n)$ to sort, plus $O(n^2)$ to fill
the table: $O(n^2)$ overall.
\textbf{Space Complexity.} $O(n)$ for \textit{dp} and \textit{parent}.
\end{complexitybox}

\begin{takeawaybox}
Whenever a pairwise condition is \textbf{transitive} (if $A$ relates to
$B$ and $B$ relates to $C$ then $A$ relates to $C$), sorting first and
then checking only consecutive chain links is enough -- the exact
property that lets both LIS's ``less than'' and this problem's
``divides'' collapse an all-pairs condition down to a simple
left-to-right ending-here DP.
\end{takeawaybox}

\section{Longest String Chain}
\label{sec:dp-longest-string-chain}

\problemstatement{Given a list of words, a \textbf{string chain} is a
sequence $\textit{word}_1, \textit{word}_2, \ldots, \textit{word}_k$
where each $\textit{word}_i$ (for $i>1$) can be formed by inserting
exactly one character, anywhere, into $\textit{word}_{i-1}$. Find the
length of the longest possible chain using the given words.}

\begin{patternbox}
A third instance of Section~\ref{sec:dp-lis-intro}'s ending-here
template: sort by \textbf{length} (the natural analogue of LIS's
numeric ordering), and replace the extension condition with ``removing
exactly one character from $\textit{word}_i$ yields exactly
$\textit{word}_j$.''
\end{patternbox}

\subsection*{The predecessor check}

A word $w_j$ is a valid predecessor of $w_i$ exactly when
$|w_i| = |w_j| + 1$ and deleting some single character from $w_i$
produces $w_j$ exactly. Checking this directly: try deleting each
character of $w_i$ in turn (one position at a time) and see if any
resulting string equals $w_j$.

\subsection*{The recurrence}

After sorting words by increasing length:
\[
  \textit{dp}[i] = 1 + \max_{j<i,\ \textit{isPredecessor}(w_j, w_i)}
  \textit{dp}[j] \quad (\text{or } 1 \text{ if no such } j).
\]

\subsection*{C++ implementation}

\begin{lstlisting}
#include <bits/stdc++.h>
using namespace std;

bool isPredecessor(string& shorter, string& longer) {
    if (longer.size() != shorter.size() + 1) return false;
    int i = 0, j = 0;
    bool skippedOne = false;
    while (i < (int)longer.size() && j < (int)shorter.size()) {
        if (longer[i] == shorter[j]) {
            i++; j++;
        } else if (!skippedOne) {
            i++; skippedOne = true;
        } else {
            return false;
        }
    }
    return true;
}

int longestStringChain(vector<string> words) {
    sort(words.begin(), words.end(),
         [](const string& a, const string& b) { return a.size() < b.size(); });

    int n = words.size();
    vector<int> dp(n, 1);
    int best = 1;

    for (int i = 0; i < n; i++) {
        for (int j = 0; j < i; j++) {
            if (isPredecessor(words[j], words[i]) && 1 + dp[j] > dp[i]) {
                dp[i] = 1 + dp[j];
            }
        }
        best = max(best, dp[i]);
    }
    return best;
}

int main() {
    vector<string> words = {"a", "b", "ba", "bca", "bda", "bdca"};
    cout << longestStringChain(words) << endl; // 4
    return 0;
}
\end{lstlisting}

\subsection*{Dry run}

\dryrun{Take $\textit{words} = [\text{"a"}, \text{"b"}, \text{"ba"},
\text{"bca"}, \text{"bda"}, \text{"bdca"}]$, already sorted by length.}

$\textit{dp}[\text{"a"}]=\textit{dp}[\text{"b"}]=1$.
$\textit{dp}[\text{"ba"}]=2$ (predecessor \texttt{"a"} or \texttt{"b"},
either gives $1+1=2$). $\textit{dp}[\text{"bca"}]=3$ (predecessor
\texttt{"ba"}, delete \texttt{'c'}). $\textit{dp}[\text{"bda"}]=3$
(predecessor \texttt{"ba"}, delete \texttt{'d'}).
$\textit{dp}[\text{"bdca"}]=4$ (predecessor \texttt{"bca"}, delete
\texttt{'d'} -- or equally \texttt{"bda"}, delete \texttt{'c'}; either
extends a length-$3$ chain to $4$). The maximum is $4$, realized by the
chain $\text{"a"} \to \text{"ba"} \to \text{"bca"} \to \text{"bdca"}$.

\begin{complexitybox}
\textbf{Time Complexity.} $O(n \log n)$ to sort, plus $O(n^2 \cdot L)$
to fill the table ($L$ = max word length, the cost of each
\textit{isPredecessor} check): $O(n^2 \cdot L)$ overall.
\textbf{Space Complexity.} $O(n)$ for \textit{dp}.
\end{complexitybox}

\begin{takeawaybox}
The ``sort, then ending-here chain with a custom extension check''
template now has three independent instances (numeric order for LIS,
divisibility for the previous section, and length-plus-single-insertion
here) -- a strong signal that this template, not any one specific
condition, is the reusable idea worth carrying forward.
\end{takeawaybox}

\section{Longest Bitonic Subsequence}
\label{sec:dp-longest-bitonic-subsequence}

\problemstatement{Given an array $\textit{arr}$, find the length of the
longest \textbf{bitonic} subsequence: one that first strictly increases,
then strictly decreases (either the increasing or decreasing part may
be empty, but not both).}

\begin{patternbox}
\emph{``Increases then decreases''} is solved by running
Section~\ref{sec:dp-lis}'s ending-here LIS computation \textbf{twice} --
once forward (for the increasing run up to each index) and once
backward (for the decreasing run from each index onward) -- then
combining the two values at every index, since a bitonic sequence's
peak can be at any position.
\end{patternbox}

\subsection*{Deriving the combination formula}

Let $\textit{inc}[i]$ = length of the longest increasing subsequence
\emph{ending} at index $i$ (exactly Section~\ref{sec:dp-lis}'s table).
Let $\textit{dec}[i]$ = length of the longest decreasing subsequence
\emph{starting} at index $i$ -- computed identically, but scanning
\emph{backward} from the end and requiring $\textit{arr}[i] >
\textit{arr}[j]$ for $j>i$ (equivalently, LIS run on the array's mirror
image). If index $i$ is the peak of a bitonic sequence, the increasing
run contributes $\textit{inc}[i]$ elements and the decreasing run
contributes $\textit{dec}[i]$ elements, but $\textit{arr}[i]$ itself is
counted in both, so:
\[
  \textit{bitonic}[i] = \textit{inc}[i] + \textit{dec}[i] - 1.
\]
The answer is $\max_i \textit{bitonic}[i]$.

\subsection*{C++ implementation}

\begin{lstlisting}
#include <bits/stdc++.h>
using namespace std;

int longestBitonicSubsequence(vector<int>& arr) {
    int n = arr.size();
    vector<int> inc(n, 1), dec(n, 1);

    for (int i = 0; i < n; i++) {
        for (int j = 0; j < i; j++) {
            if (arr[j] < arr[i]) {
                inc[i] = max(inc[i], 1 + inc[j]);
            }
        }
    }

    for (int i = n - 1; i >= 0; i--) {
        for (int j = n - 1; j > i; j--) {
            if (arr[j] < arr[i]) {
                dec[i] = max(dec[i], 1 + dec[j]);
            }
        }
    }

    int best = 0;
    for (int i = 0; i < n; i++) {
        best = max(best, inc[i] + dec[i] - 1);
    }
    return best;
}

int main() {
    vector<int> arr = {1, 11, 2, 10, 4, 5, 2, 1};
    cout << longestBitonicSubsequence(arr) << endl; // 6
    return 0;
}
\end{lstlisting}

\subsection*{Dry run}

\dryrun{Take $\textit{arr} = [1,11,2,10,4,5,2,1]$.}

$\textit{inc} = [1,2,2,3,3,4,2,1]$ (the run $1,2,4,5$ ending at index
$5$ gives $\textit{inc}[5]=4$). $\textit{dec} = [1,5,2,4,3,3,2,1]$ (the
run $11,10,5,2,1$ starting at index $1$ gives $\textit{dec}[1]=5$).
$\textit{bitonic}[1] = \textit{inc}[1]+\textit{dec}[1]-1 = 2+5-1=6$,
the maximum -- realized by the sequence $1,11,10,5,2,1$, increasing
from $1$ to the peak $11$ (index $1$) and then decreasing.

\begin{complexitybox}
\textbf{Time Complexity.} $O(n^2)$ -- two nested-loop passes.
\textbf{Space Complexity.} $O(n)$ for \textit{inc} and \textit{dec}.
\end{complexitybox}

\begin{takeawaybox}
When a problem's structure splits naturally into a ``before the
turning point'' half and an ``after the turning point'' half, compute
each half's ending-here table independently (one forward, one
backward) and combine them at every candidate turning point -- a
pattern that generalizes well beyond bitonic sequences to any
problem with a single peak or pivot.
\end{takeawaybox}

\section{Number of Longest Increasing Subsequences}
\label{sec:dp-number-of-lis}

\problemstatement{Given $\textit{arr}$, find the \textbf{number of
distinct} longest increasing subsequences (not just their shared
maximum length).}

\begin{patternbox}
This closing problem of the chapter pairs Section~\ref{sec:dp-lis}'s
ending-here \textit{dp} table with a second, parallel table
$\textit{count}[i]$: \textbf{how many} subsequences of the maximum
length achieve $\textit{dp}[i]$ -- updated using the same
``exists / count'' duality first introduced in the DP on Subsequences
chapter (there distinguishing whether a target sum is
\emph{achievable} versus \emph{how many ways} achieve it).
\end{patternbox}

\subsection*{Deriving the count update rule}

While filling $\textit{dp}[i]$ in the usual way, also track
$\textit{count}[i]$, initialized to $1$ (every index is trivially one
subsequence of length $1$, by itself). For each $j<i$ with
$\textit{arr}[j]<\textit{arr}[i]$:
\begin{itemize}
  \item If $\textit{dp}[j]+1 > \textit{dp}[i]$: a \textbf{strictly
        longer} chain through $j$ has just been discovered. Update
        $\textit{dp}[i] \gets \textit{dp}[j]+1$ and \textbf{reset}
        $\textit{count}[i] \gets \textit{count}[j]$ (every previous,
        shorter-chain count at $i$ is now obsolete -- it counted chains
        of a length that's no longer the best).
  \item If $\textit{dp}[j]+1 = \textit{dp}[i]$: this is \textbf{another
        equally good} way to reach the current best length at $i$.
        \textbf{Accumulate}: $\textit{count}[i] \mathrel{+}=
        \textit{count}[j]$.
\end{itemize}
After filling both tables, let $\textit{maxLen} = \max_i \textit{dp}[i]$.
The answer is $\sum_{i:\ \textit{dp}[i]=\textit{maxLen}}
\textit{count}[i]$ (every index whose best chain achieves the global
maximum length contributes its count).

\subsection*{C++ implementation}

\begin{lstlisting}
#include <bits/stdc++.h>
using namespace std;

int numberOfLIS(vector<int>& arr) {
    int n = arr.size();
    vector<int> dp(n, 1), count(n, 1);

    for (int i = 0; i < n; i++) {
        for (int j = 0; j < i; j++) {
            if (arr[j] < arr[i]) {
                if (dp[j] + 1 > dp[i]) {
                    dp[i] = dp[j] + 1;
                    count[i] = count[j];
                } else if (dp[j] + 1 == dp[i]) {
                    count[i] += count[j];
                }
            }
        }
    }

    int maxLen = *max_element(dp.begin(), dp.end());
    int total = 0;
    for (int i = 0; i < n; i++) {
        if (dp[i] == maxLen) total += count[i];
    }
    return total;
}

int main() {
    vector<int> arr = {1, 3, 5, 4, 7};
    cout << numberOfLIS(arr) << endl; // 2
    return 0;
}
\end{lstlisting}

\subsection*{Dry run}

\dryrun{Take $\textit{arr} = [1,3,5,4,7]$.}

\begin{center}
\begin{tabular}{c|c|c|c|c|c}
$i$ & 0 & 1 & 2 & 3 & 4 \\ \hline
$\textit{arr}[i]$ & 1 & 3 & 5 & 4 & 7 \\ \hline
$\textit{dp}[i]$ & 1 & 2 & 3 & 3 & 4 \\ \hline
$\textit{count}[i]$ & 1 & 1 & 1 & 1 & 2
\end{tabular}
\end{center}

At $i=3$ (value $4$): only $j=0,1$ qualify ($1<4$, $3<4$); the best is
through $j=1$ ($\textit{dp}[1]+1=3$), giving $\textit{dp}[3]=3$,
$\textit{count}[3]=\textit{count}[1]=1$. At $i=4$ (value $7$):
$j=2$ gives $\textit{dp}[2]+1=4$ (new max, reset
$\textit{count}[4]=\textit{count}[2]=1$); then $j=3$ \emph{also} gives
$\textit{dp}[3]+1=4$ (ties the current best, accumulate:
$\textit{count}[4] \mathrel{+}= \textit{count}[3] = 1+1=2$).
$\textit{maxLen}=4$, achieved only at $i=4$, contributing
$\textit{count}[4]=2$: the two LIS of length $4$ are $1,3,5,7$ and
$1,3,4,7$.

\begin{complexitybox}
\textbf{Time Complexity.} $O(n^2)$.
\textbf{Space Complexity.} $O(n)$ for \textit{dp} and \textit{count}.
\end{complexitybox}

\begin{takeawaybox}
This chapter closes exactly where the DP on Subsequences chapter's
``exists / count / optimize'' framing predicted it would: once a
length-optimizing ending-here table exists, a parallel counting table
updated by the same reset-on-strict-improvement,
accumulate-on-tie rule answers ``how many'' for free. The same two-table
pairing pattern is worth trying on any optimization DP whenever a
follow-up problem asks for a count instead of (or in addition to) the
optimum itself.
\end{takeawaybox}

